\section{Гидростатика}
%%%%%%%%%%%%%%%%%%%%%%%%%%%%%%%%%%%%%%%%%%%%%%%%%%%%
\subsection{Давление}
%%%%%%%%%%%%%%%%%%%%%%%%%%%%%%%%%%%%%%%%%%%%%%%%%%%%
%%%%%%%%%%%%%%%%%%%%%%%%%%%%%%%%%%%%%%%%%%%%%%%%%%%%
\z{Шкаф массой $100\,\text{кг}$ стоит на четырёх ножках, след каждой из которых имеет форму квадрата. Чему равна сторона этого квадрата, если шкаф оказывает давление на пол $150\,\text{кПа}$?}
%
{$a = 4\,\text{см}$
}

%%%%%%%%%%%%%%%%%%%%%%%%%%%%%%%%%%%%%%%%%%%%%%%%%%%%
\z{Плоскодонная баржа получила пробоину в днище площадью $200\,\text{см}^2$. С какой силой надо прижимать пластырь, которым заделывают пробоину, чтобы выдержать напор воды на глубине $2\,\text{м}$?}
%
{$F = 400\,\text{Н}$
}

%%%%%%%%%%%%%%%%%%%%%%%%%%%%%%%%%%%%%%%%%%%%%%%%%%%%
\z{Определите давление воды на дно сосуда, если сосуд полностью заполнен водой и имеет форму куба, а масса воды в нём равна $64\,\text{г}$.
}
%
{$400\,\text{Па}$
}

%%%%%%%%%%%%%%%%%%%%%%%%%%%%%%%%%%%%%%%%%%%%%%%%%%%%
\z{На дно аквариума длиной $40\,\text{см}$ и шириной $25\,\text{см}$ положили чугунный шар массой $700\,\text{г}$. На сколько увеличилось после этого давление воды на дно, если вода из аквариума не вылилась? Шар погружён в воду полностью.
}
{$\Delta P = \frac{\rho_{\text{в}} g m}{\rho_{\text{св}} S} = 10\,\text{Па}$
}

%%%%%%%%%%%%%%%%%%%%%%%%%%%%%%%%%%%%%%%%%%%%%%%%%%%%
\z{До какой высоты $h$ нужно налить жидкость в цилиндрическое ведро радиуса $R$, чтобы сила $F$, с которой жидкость давит на боковую поверхность сосуда, была равна силе давления на дно?
}
%
{$h = r$
}

%%%%%%%%%%%%%%%%%%%%%%%%%%%%%%%%%%%%%%%%%%%%%%%%%%%%
\z{Брусок массой $m = 2\,\text{кг}$ имеет форму параллелепипеда. Лёжа на одной из граней, он оказывает давление $p_1 = 1\,\text{кПа}$, лёжа на другой — давление $p_2 = 2\,\text{кПа}$, стоя на третьей — давление $p_3 = 4\,\text{кПа}$. Каковы размеры бруска?
}
%
{$5\,\text{см},\ 10\,\text{см},\ 20\,\text{см}$
}


%%%%%%%%%%%%%%%%%%%%%%%%%%%%%%%%%%%%%%%%%%%%%%%%%%%%
\z{На горизонтальном листе резины лежит перевёрнутая кастрюля радиусом $R = 10\,\text{см}$ и высотой $H = 15\,\text{см}$. В дне кастрюли просверлено круглое отверстие радиуса $r = 1\,\text{см}$, в которое плотно вставлена лёгкая вертикальная трубка. В кастрюлю через трубку наливают воду. Когда вода заполняет всю кастрюлю и поднимается по трубке на $h = 4\,\text{см}$, она начинает вытекать из-под краёв кастрюли. Какова масса $m$ кастрюли?
}
%
{$m = h \left( R^2 - r^2 \right) \rho = 1.24\,\text{кг}$
}

